\documentclass[11pt]{article}
\usepackage[top=0.85in,bottom=0.85in,left=1in,right=1in]{geometry}
\usepackage{times}
\usepackage{parskip}
\setlength{\parskip}{6pt}
\usepackage{hyperref}
\hypersetup{colorlinks=true, urlcolor=blue, linkcolor=black}
\pagestyle{empty}

\begin{document}

\noindent
May 24, 2026

\bigskip

\noindent
The Editors\\
\textit{Nature}

\bigskip

\noindent
Dear Editors,

\medskip

We submit for your consideration ``Selection as Learning: A Formal Unification of
Evolutionary Quantitative Genetics and PAC Learning Theory.''

The central claim is a theorem, not an analogy: Fisher's Fundamental Theorem of
Natural Selection, the Price Equation, and PAC generalization bounds are
instances of empirical risk minimization on a Fisher-information manifold under a
capacity constraint. This shared geometric structure means that natural selection
\emph{is} natural gradient ascent---a connection that resolves longstanding
interpretive debates about the FTNS, extends Valiant's PAC-evolvability
framework from binary to continuous polygenic phenotypes, and yields the first
quantitative bridge between population genetics and statistical learning theory.

The empirical value is immediate. Because the unified framework produces a
computable threshold---the evolutionary sample complexity $N_e^*$, the tumour
cell burden below which reliable resistance evolution is mathematically
precluded---it is testable against cancer evolutionary genomics data. Across nine
cancer types, the theoretical PAC bound $\varepsilon^*(N_e, V_A)$ correlates with
observed resistance durability at $r = 0.96$ ($p < 0.0001$). Across twenty cancer
types, the FTNS convergence rate $V_A/\bar{W}$ correlates with measured selection
coefficients at $r = 0.999$ ($p < 0.0001$). Both $N_e$ and $V_A$ are already
delivered by standard clinical whole-genome sequencing pipelines. The threshold
$N_e^*$ is computable at treatment initiation.

The clinical consequence is equally direct. Standard-of-care sequential
monotherapy inadvertently provides tumours with the iterative training loops that
maximise evolutionary learning: each drug line is one labelled training example,
each relapse a weight update. Upfront combination therapy engineered to maintain
$N_e < N_e^*$ closes the evolutionary search space entirely before resistance can
be reliably discovered. The mathematics does not suggest sequential monotherapy
is suboptimal. It proves it.

We believe this paper is well suited to \textit{Nature} for three reasons. First,
the theoretical result is foundational: it places 3.8 billion years of natural
selection within the formal framework of machine learning, with consequences for
evolutionary biology, statistical learning theory, and ecology. Second, the
empirical validation is immediate and reproducible from publicly available data
(cBioPortal, PCAWG, TRACERx 421), with all code deposited at
\href{https://github.com/rstil2/selection-as-learning}{github.com/rstil2/selection-as-learning}.
Third, the clinical application---a computable resistance threshold derivable at
the bedside from existing sequencing data---speaks directly to one of the most
consequential unsolved problems in oncology.

The manuscript has not been published or submitted elsewhere. We have no
competing interests to declare.

\bigskip

\noindent
Sincerely,

\bigskip
\bigskip

\noindent
R.\ Craig Stillwell\\
Independent Researcher\\
\href{mailto:craig.stillwell@gmail.com}{craig.stillwell@gmail.com}

\end{document}
